% =============================================================================
% Anhang: Vollständiges Findings-Register
% Arbeitsstand: 03.09.2026
%
% Zweck:
% Vollständige und transparente Dokumentation der formativen WP-12 Findings.
% Der Haupttext, insbesondere Kapitel 7, verwendet nur eine verdichtete Auswahl.
%
% Statuslogik:
% - Historische Zustände werden NICHT überschrieben.
% - Ein späterer Retest wird als finaler / späterer Zustand separat dokumentiert.
% - "Open" bedeutet nicht automatisch thesis-kritisch.
% - OBS-Findings sind überwiegend UX-/Integrationsbeobachtungen und keine
%   wissenschaftlichen Kernergebnisse der Arbeit.
%
% Primäre Projektevidenz:
% collaboration/audits/wp12_findings.md
% collaboration/audits/wp12_formative_self_test_report.md
% collaboration/audits/post_golden_finding_triage.md
% collaboration/roadmap.md
% relevante accepted checkpoints / ADRs
%
% WICHTIG:
% Dieses Register ist thesisgerecht verdichtet. Die technischen Primärartefakte
% und vollständigen Detailtexte verbleiben im Repository.
% =============================================================================

\section{Vollständiges Findings-Register}
\label{app:findings_register}

\subsection{Leselogik und Statusführung}

Die formative Verifikation des Prototyps erzeugte technische Blocker, semantische
Findings, UX- und Integrationsbeobachtungen sowie positive Verifikationsergebnisse.
Die vollständige Auflistung wird im Anhang erhalten, damit auch behobene oder für
den Haupttext nachrangige Befunde transparent nachvollziehbar bleiben. Historische
Fehlerzustände werden dabei nicht durch spätere erfolgreiche Retests ersetzt.
Stattdessen werden Erstbeobachtung und finaler Status gemeinsam dokumentiert.

Die verwendeten Präfixe haben folgende Bedeutung:

\begin{itemize}
    \item \texttt{BLK}: blockierendes Finding in einem verifizierten Workflow- oder Gate-Pfad,
    \item \texttt{SEM}: semantisches oder modellierungsbezogenes Finding,
    \item \texttt{OBS}: formative UX-, Workflow- oder Integrationsbeobachtung,
    \item \texttt{PASS}: positive oder bestätigende Projektevidenz.
\end{itemize}

\subsection{Blockierende Findings}

\paragraph{BLK-001: Derivation Producer Contract zu permissiv}

\textbf{Erstbeobachtung:}
Der Output-Contract des Derivation-Assessors ließ formal plausible, aber für den
strikt validierenden Downstream-Adapter nicht zuverlässig adaptierbare Outputs zu.

\textbf{Reaktion:}
Härtung des Producer Contracts. Semantische Unsicherheit wurde von
Identitäts-, Provenance- und strukturellen Integritätsfehlern getrennt.

\textbf{Finaler Status:}
\texttt{CORRECTED / focused validation passed}.

\textbf{Thesis-Verwendung:}
Unterstützende formative Evidenz für Contract- und Fail-Closed-Grenzen, kein
eigenständiges Kernergebnis in Kapitel 7.


\paragraph{BLK-002: Cross-Source Processing Artifact Identity Collision}

\textbf{Historische Erstbeobachtung:}
Der formale Multi-Source-Pfad konnte Processing-Artifacts projektweit nicht
ausreichend eindeutig an Source-, Run- und Artifact-Kontext binden.
Der Stage-A-Test wurde deshalb als \texttt{FAILED WITH BLOCKER} klassifiziert.

\textbf{Korrektur:}
Kontextuelle Processing-Artifact-Identität, exakter Project Fit,
\texttt{ProjectFitPhaseHHandoff}, source-scoped Phase-H Subject Identity,
getrennte source-lokale AEI-Sets und ein gemeinsamer Model-Candidate-Pfad.

\textbf{Retest:}
Project \texttt{308131} durchlief den Genuine-Multi-Source-Pfad erfolgreich.

\textbf{Finaler Status:}
\texttt{RESOLVED / LIVE RETESTED}. Kein WP-12-Blocker verbleibt aus BLK-002.

\textbf{Verbleibende Scope-Grenze:}
Automatische projektweite semantische Reconciliation und Change Control bleiben
Ausblick. Concern-centric S3A/S3B/S4/S5 bleiben Prototype Evidence.

\textbf{Thesis-Verwendung:}
Kernergebnis für Multi-Source Identity, Provenance und Project Fit.


\paragraph{BLK-003: Semantic Effectiveness / Engineering-Subject Quality}

\textbf{Historische Erstbeobachtung:}
Die frühere Persona-first-Architektur erzeugte Engineering Subjects zu spät und
zu stark aus agentengenerierten Namen, Typen und Relationship-Formulierungen.
Processing- und Orchestration-Kontext konnte zusätzlich die semantische Population
verunreinigen.

\textbf{Architekturkorrektur:}
R4c mit source-grounded Evidence Detection, gemeinsamer Subject Discovery,
Canonical Subjects vor Persona Interpretation und expliziter Human Authority.

\textbf{Retest:}
Der Single-Source-Live-Pfad auf Project \texttt{120412} erreichte den
source-grounded Human-Review- und später den vollständigen Golden-E2E-Pfad.

\textbf{Finaler Status:}
\texttt{RESOLVED / single-source E2E validated}.

\textbf{Thesis-Verwendung:}
Kernergebnis für Source Grounding und Subject Identity vor Persona Interpretation.


\paragraph{BLK-004: Approved Input Promotion fails for R4c semantic classification}

\textbf{Erstbeobachtung:}
Die R4c-Klassifikation war als Tupel aus \texttt{information_type},
\texttt{modality} und \texttt{epistemic_status} modelliert, während der
Legacy-Promotion-Pfad einen einzelnen skalaren Klassifikationswert erwartete.

\textbf{Retest Project 120412:}
17 Approved Inputs wurden erzeugt, sieben Open Questions bewusst übersprungen.
Human Overrides blieben erhalten.

\textbf{Finaler Status:}
\texttt{CORRECTED / LIVE VALIDATED}.

\textbf{Thesis-Verwendung:}
Unterstützende Contract-Evidenz, kein eigenständiges Kernergebnis.


\paragraph{BLK-005: Approved Engineering Information not authoritative Phase-H input}

\textbf{Erstbeobachtung:}
Der Phase-H-Handoff verwendete zunächst Approved Inputs, aber nicht die vollständige
human-autorisierte Engineering Information einschließlich akzeptierter semantischer
Relationships.

\textbf{Korrektur:}
Der Phase-H-Request bindet Approved Subjects und Accepted Semantic Relationships
als autoritative AEI-Population. Relationships werden nicht künstlich in
Legacy-Approved-Input-Formate umgedeutet.

\textbf{Retest Project 120412:}
17 Approved Subjects und 21 akzeptierte semantische Relationships wurden
gebunden. Sechs Relationships zu nicht model-promotierbaren Open Questions wurden
explizit als \texttt{intentionally\_not\_projected} geführt.

\textbf{Finaler Status:}
\texttt{CORRECTED / LIVE VALIDATED}.

\textbf{Thesis-Verwendung:}
Kernergebnis zur Vollständigkeit des autoritativen Engineering-Information-Handoffs.


\paragraph{BLK-006: LLM-assisted Model Proposal generation fails in live E2E}

\textbf{Historische Erstbeobachtung:}
Der erste echte LLM-assisted Model-Proposal-Pfad des R4c-Live-E2E blockierte den
Downstream-Modellpfad. \texttt{IEM-000001} enthielt 13 Elemente und drei
Relationships; der Generation-Preflight meldete vier nicht unterstützte Mappings.
Die Dependency-Relationship war unterstützt.

\textbf{Architekturkorrektur:}
Die Verantwortung zwischen Engineering Meaning, Target-Model Representation,
Target-Model Formulation und deterministischer Serialization wurde expliziter
getrennt. Der Golden-E2E-Pfad wurde anschließend mit der korrigierten Architektur
durchlaufen.

\textbf{Finaler Status:}
\texttt{RESOLVED BY ARCHITECTURAL CORRECTION / Golden E2E retest PASS}.

\textbf{Thesis-Verwendung:}
Kernergebnis für die Trennung von Meaning, Representation, Formulation und
Serialization.


\paragraph{BLK-007: Model Assembly invoked Model Placement personas for Relationship representation}

\textbf{Erstbeobachtung:}
Model Assembly delegierte Relationship-Repräsentation an die Placement-Personas
und überschritt damit die akzeptierte Verantwortungsgrenze.

\textbf{Korrektur:}
Model Assembly wurde auf deterministische Materialisierung autorisierter Semantik
begrenzt. Nicht exakt abbildbare Semantik bleibt unresolved und wird nicht durch
eine erneute LLM-Interpretation während der Assembly ersetzt.

\textbf{Retest Project 120412:}
Assembly Draft mit 13 Elementen und drei Relationships, ohne Relationship-Variance;
drei Zielrepräsentationen blieben explizit unresolved/unmapped.

\textbf{Finaler Status:}
\texttt{CORRECTED / LIVE RETEST PASS}.

\textbf{Thesis-Verwendung:}
Kernergebnis zur deterministischen Assembly-Grenze.


\paragraph{WP12-BLK-SEM015-L-001: Final Review Handoff}

\textbf{Erstbeobachtung:}
Der finale Authority-Handoff zwischen akzeptierter Final-Review-Revision,
Human Release und Published Output musste als geschlossene Kette verifiziert werden.

\textbf{Retest / Ergebnis:}
\texttt{FRV-000002 -> FRD-000001 -> OUT-000001}.

\textbf{Finaler Status:}
\texttt{RESOLVED}.

\textbf{Thesis-Verwendung:}
Unterstützende Evidenz für die Trennung technischer Validation, Human Review und
Publication.


\subsection{Semantische und modellierungsbezogene Findings}

\paragraph{SEM-001: Source Purity Regression}

\textbf{Beobachtung:}
Orchestrierungs-, Task- und Prozessinformationen wurden teilweise als positive
Engineering Information interpretiert.

\textbf{Korrektur:}
Positive Engineering Information wird an source-grounded Content gebunden.

\textbf{Status:}
\texttt{CORRECTED IN R4c / LIVE SINGLE-SOURCE VALIDATED}. Eine separate,
dedizierte Revalidierung genau dieses Findings im Multi-Source-Test ist nicht
als eigener Finding-Retest dokumentiert.

\textbf{Thesis-Verwendung:}
Kernergebnis Source Boundary / Source Purity.


\paragraph{SEM-002: Model Relevance Regression}

\textbf{Beobachtung:}
Semantisch plausible Extraktion führte nicht automatisch zu einer für die
Modellableitung geeigneten Repräsentation.

\textbf{Recovery:}
Subject-level Recovery wurde live validiert. Der Downstream-Modellpfad wurde
später im Golden E2E validiert; allgemeine Construct-Coverage bleibt über
SEM-011 begrenzt.

\textbf{Thesis-Verwendung:}
Kernergebnis zur Trennung von Engineering Meaning und Zielmodell-Repräsentation.


\paragraph{SEM-003: Persona-driven Subject Multiplication}

\textbf{Beobachtung:}
Personas erzeugten weitgehend unabhängige Proposal-Populationen. Engineering
Subject Identity entstand dadurch zu spät.

\textbf{Korrektur:}
Gemeinsame Evidence- und Subject Identity wird vor Persona Interpretation
etabliert.

\textbf{Status:}
\texttt{CORRECTED ARCHITECTURALLY / LIVE VALIDATED}.

\textbf{Thesis-Verwendung:}
Kernergebnis.


\paragraph{SEM-004: Cross-Unit Consolidation empirisch ineffektiv}

\textbf{Beobachtung:}
Im Run \texttt{877791 / RUN-000001} lagen nach D3 109 Subjects und nach D4
weiterhin 109 synthetisierte Subjects vor.

\textbf{Status:}
\texttt{SUPERSEDED BY R4c ARCHITECTURE / HISTORICAL FINDING}.

\textbf{Thesis-Verwendung:}
Historische quantitative Evidenz für Kapitel 7. Die Interpretation der
nachgelagerten Konsolidierung gehört in Kapitel 8.


\paragraph{SEM-005: Review-Item Count ist kein Optimierungsziel}

\textbf{Beobachtung:}
Die reine Zahl der Review Items ließ keine belastbare Aussage über
Source Purity, Grounding, Relevance oder tatsächliche Persona-Varianz zu.

\textbf{Status:}
\texttt{ACCEPTED INTERPRETATION OF TEST RESULT}.

\textbf{Thesis-Verwendung:}
Unterstützung für die Diskussion von Qualitäts- und Effizienzmetriken.


\paragraph{SEM-006: Human Review Cards zu textlastig}

\textbf{Beobachtung:}
Review Cards erforderten zu viel Rationale- und Detaillektüre, bevor die
eigentliche Engineering-Aussage sichtbar wurde.

\textbf{Status:}
\texttt{OPEN / UX + ENGINEERING PRESENTATION} beziehungsweise ohne dokumentierten
abschließenden Revalidation-Status.

\textbf{Thesis-Verwendung:}
Anhang / UX-Limitation.


\paragraph{SEM-007: Source-wide Completeness Request skaliert schlecht}

\textbf{Beobachtung:}
Project \texttt{691616} erzeugte bei drei Personas und zwei Runs 198 erfolgreiche
Stage-01..03-LLM-Responses. Die globale Completeness-Anfrage wurde praktisch zu groß.

\textbf{Status:}
\texttt{HISTORICAL / PRODUCTIVE R4c PATH SUPERSEDES SOURCE-WIDE COMPLETENESS}.

\textbf{Thesis-Verwendung:}
Unterstützende Evidenz für Laufzeit, Call-Volumen und bounded work units in
Kapitel 8.


\paragraph{SEM-008: Legitimate uncertainty funktioniert als Human-Review-Fall}

\textbf{Beobachtung:}
Im Run \texttt{877791} entstand eine legitime Open Question zur Klassifikation
von \texttt{microscope workstation}.

\textbf{Status:}
\texttt{POSITIVE FINDING}.

\textbf{Thesis-Verwendung:}
Kernergebnis für Unsicherheit und Human Authority.


\paragraph{SEM-009: Predicate-variant Relationship Consolidation unvollständig}

\textbf{Beobachtung:}
Semantisch nahe oder äquivalente Relationship-Hypothesen können bei
unterschiedlichen Prädikaten getrennt bleiben.

\textbf{Status:}
\texttt{OPEN / NON-BLOCKING SEMANTIC QUALITY}.

\textbf{Thesis-Verwendung:}
PoC-Grenze / Anhang, gegebenenfalls Kapitel 7.7.


\paragraph{SEM-010: Relationship lifecycle nicht automatisch an rejected Subject ausgerichtet}

\textbf{Beobachtung:}
Eine akzeptierte Relationship zu einem später abgelehnten Subject blieb zunächst
accepted und blockierte die Finalization, bis sie explizit reopened und abgelehnt
wurde.

\textbf{Status:}
\texttt{OPEN / NON-BLOCKING GOVERNANCE + SEMANTIC LIFECYCLE}.

\textbf{Thesis-Verwendung:}
PoC-Grenze und unterstützende Human-Authority-Evidenz.


\paragraph{SEM-011: Incomplete SysML element-type coverage}

\textbf{Beobachtung:}
Die aktuelle Taxonomie und der Target-Pfad decken nicht alle relevanten
SysML-v2-Constructs ab, unter anderem spezifische Ports, Interfaces, Parts,
Actions/Functions, States oder Occurrences.

\textbf{Status:}
\texttt{OPEN / ACCEPTED MVP LIMITATION / NON-BLOCKING}.

\textbf{Thesis-Verwendung:}
Zentrale Claim Boundary des PoC.


\paragraph{SEM-012: Engineering Information Type und Target Model Representation nicht ausreichend getrennt}

\textbf{Beobachtung:}
Engineering Classification beschreibt Bedeutung, schreibt aber nicht eindeutig
den Zielmodell-Construct vor. Im Live-Pfad konnte beispielsweise Engineering
Information vom Typ Constraint als Stakeholder Requirement repräsentiert werden.

\textbf{Finaler Status:}
\texttt{IMPLEMENTED IN CURRENT SINGLE-SOURCE PATH} für den validierten Scope.

\textbf{Thesis-Verwendung:}
Kernergebnis.


\paragraph{SEM-013: Shared Placement Ambiguity wird mit Persona Variance vermischt}

\textbf{Beobachtung:}
Wenn alle Placement-Personas dieselbe Menge mehrerer Alternativen liefern, liegt
gemeinsame Ambiguität und nicht zwingend Inter-Persona-Varianz vor.

\textbf{Status:}
\texttt{ACCEPTED FINDING / OPEN}.

\textbf{Thesis-Verwendung:}
Kernergebnis für TF2, zusammen mit Human-Interaction-Limitationen.


\paragraph{SEM-014: Unresolved Target Relationship Representation benötigt explizite Human Selection}

\textbf{Beobachtung:}
Ein unresolved Relationship-Fall darf nicht über einen inhaltlichen UI-Default
implizit Engineering Authority erhalten.

\textbf{Status:}
\texttt{ACCEPTED FINDING / OPEN}.

\textbf{Thesis-Verwendung:}
Kernergebnis Human Authority / Target Representation.


\paragraph{SEM-015: Target-Model Formulation als explizite Processing Stage}

\textbf{Historische Beobachtung:}
Der Pfad ging zu direkt von approved Engineering Meaning und Placement in Richtung
deterministischer SysML-v2-Serialization.

\textbf{Korrektur:}
Explizite Trennung von Engineering Meaning, Target-Model Representation,
Target-Model Formulation und Serialization.

\textbf{Finaler Status:}
Architektur implementiert, Golden-E2E-Scope validiert, allgemeine
Target-Type-Formulation weiterhin \texttt{PARTIAL / OPEN}. Die Generalisierung ist
an SEM-011 gekoppelt.

\textbf{Thesis-Verwendung:}
Kernergebnis.


\paragraph{SEM-015-F01: Dependency-aware in-review regeneration}

\textbf{Beobachtung:}
Eine geänderte Klassifikation erfordert im aktuellen PoC einen neuen fokussierten
Refinement-Run statt selektiver automatischer Invalidierung und Regeneration.

\textbf{Status:}
\texttt{DEFERRED / NOT REQUIRED FOR CURRENT PoC}.

\textbf{Thesis-Verwendung:}
Ausblick / Performance- und Workflowoptimierung.


\subsection{Formative UX- und Integrationsbeobachtungen}

% Die OBS-Liste wird vollständig erhalten, aber nicht als wissenschaftliche
% Hauptergebnisstruktur verwendet.

\begin{description}
    \item[OBS-001] Project selection lost on toggle. Status: offen beziehungsweise nicht abschließend revalidiert.
    \item[OBS-002] Add-first-source route / discoverability. Status: offen beziehungsweise nicht abschließend revalidiert.
    \item[OBS-003] Duplicate Create Project interaction. Status: offen beziehungsweise nicht abschließend revalidiert.
    \item[OBS-004] Multi-file add affordance. Status: offen beziehungsweise nicht abschließend revalidiert.
    \item[OBS-005] Drag-and-drop first-click behavior. Status: offen beziehungsweise nicht abschließend revalidiert.
    \item[OBS-006] Row-action selection scalability. Status: offen beziehungsweise nicht abschließend revalidiert.
    \item[OBS-007] Long Processing feedback. Status: offen beziehungsweise nicht abschließend revalidiert.
    \item[OBS-008] Processing cancellation architecture. Status: deferred / hoher Lifecycle-Aufwand.
    \item[OBS-009] Optional tips / guidance. Status: niedrige UX-Priorität.
    \item[OBS-010] Agent scope label clarity. Status: offen beziehungsweise nicht abschließend revalidiert.
    \item[OBS-011] Live state inconsistent after source change. Status: offen beziehungsweise nicht abschließend revalidiert.
    \item[OBS-012] Processing coupled to active Streamlit page run. Status: offene Integrations-/Lifecycle-Grenze.
    \item[OBS-013] Review Queue zeigt vor Workspace-Materialisierung potenziell irreführend null Items. Status: offen.
    \item[OBS-014] Reviewer identity confirmation. Status: offene Governance-/UX-Beobachtung.
    \item[OBS-015] Kein expliziter Reprocessing-/Supersession-Flow für bereits processierte Sources. Status: deferred / offen.
    \item[OBS-016] Persona configuration presentation. Status: offen beziehungsweise nicht abschließend revalidiert.
    \item[OBS-017] Long-running Processing feedback / performance. Status: offene praktische Limitation.
    \item[OBS-018] Global Project Selector session-state conflict. Status: offen beziehungsweise nicht abschließend revalidiert.
    \item[OBS-019] Subject Discovery over-generation / insufficient abstention. Status: \texttt{OPEN / NON-BLOCKING QUALITY}.
    \item[OBS-020] Editing Human Review decision änderte Action Semantics zunächst nicht. Status: \texttt{CORRECTED / FOCUSED + LIVE VALIDATED}.
    \item[OBS-021] Plain Accept wurde zunächst als \texttt{accepted\_with\_modification} persistiert. Status: \texttt{CORRECTED FOR NEW DECISIONS / FOCUSED + LIVE VALIDATED}.
    \item[OBS-022] Review Queue crash bei null verbleibenden Decisions. Status: \texttt{CORRECTED / RETESTED}.
    \item[OBS-023] Relationship Review zu versteckt / Blocker-Navigation unklar. Status: \texttt{OPEN / UX}.
    \item[OBS-024] Relationship Review zu ID-zentriert. Status: \texttt{OPEN / UX}.
    \item[OBS-025] Accepted Subject Decisions konnten zunächst nicht reopened werden. Status: \texttt{CORRECTED / LIVE VALIDATED}.
    \item[OBS-026] Decision Controls blieben nach Relationship Decision aktiv. Status: \texttt{CORRECTED / LIVE VALIDATED}.
    \item[OBS-027] Kein expliziter deterministischer Export des Finalized Review Reports. Status: \texttt{OPEN / SHOULD}.
    \item[OBS-028] Cross-layer identity transition zwischen \texttt{SUBJ-*} und \texttt{RIT-*} unzureichend erklärt. Status: \texttt{OPEN / UX + TRACEABILITY EXPLANATION}.
    \item[OBS-029] Post-promotion same-render status teilweise stale. Status: \texttt{OPEN / LOW UX}.
    \item[OBS-030] Kein sichtbarer Processing State während Model Proposal Generation. Status: \texttt{OPEN / SHOULD}.
    \item[OBS-031] Human Model Placement Review erfordert zu viel Interaktion. Status: \texttt{ACCEPTED FINDING / OPEN}.
    \item[OBS-032] Review-gate naming ambiguity. Status: offener Backlog nach Gate 3.
    \item[OBS-DASH-001] Dashboard authority-state representation. Status: offener Backlog nach Gate 3.
    \item[OBS-033] Im Roadmap-Closeout als \texttt{PASS / closed for validated v0.3.0 scope after successful live retest and publication} geführt. Der verfügbare Closeout-Auszug enthält keinen separaten Finding-Titel; deshalb wird hier bewusst kein Titel rekonstruiert.
\end{description}

\subsection{Positive und bestätigende Findings}

\paragraph{PASS-001: Project Setup und Source Registration grundsätzlich funktionsfähig}
Projektanlage und separate Source-Registrierung funktionierten. Der historische
BLK-002-Vorbehalt für den Multi-Source-Pfad wurde später durch den erfolgreichen
Project-308131-Retest abgelöst.

\paragraph{PASS-002: D4 zu Human Review Routing technisch funktionsfähig}
Der historische Run \texttt{877791} erreichte den Review Workspace. Dieses Finding
bezieht sich auf die frühere Processing-Architektur und bleibt als historische
Integrationsevidenz erhalten.

\paragraph{PASS-003: Human-in-the-Loop Authority Boundary bleibt erhalten}
Agent Agreement oder Consensus erzeugt keine automatische Engineering Approval.

\paragraph{PASS-004: Semantische Unsicherheit kann reviewbar bleiben}
Nicht eindeutig lösbare Klassifikationen oder Relationship-Endpunkte können als
Open Question erhalten bleiben, statt erfunden oder pauschal als Processing Failure
behandelt zu werden.

\paragraph{PASS-005: Source-bounded LLM-Grundaufgabe qualitativ machbar}
Ein qualitativer Benchmark zeigte für die verwendete Demo-Source, dass ein
allgemeines LLM source-grounded Engineering Information identifizieren konnte.
Dieses Finding ist keine allgemeine Modell-Performance-Aussage.

\paragraph{PASS-006: R4c source-grounded semantic recovery live-demonstrable}
Project \texttt{120412} erreichte mit realem LLM einen subject-centric Human Review
mit 24 Canonical Subjects.

\paragraph{PASS-007: Subject/Relationship Review und Finalization authority-preserving}
24 Subject Decisions und 30 outgoing Relationship Decisions wurden explizit
getroffen. Finalization blieb fail-closed, bis eine inkonsistente accepted
Relationship korrigiert wurde. Finalisierte Revision: \texttt{RVR-000062}.

\paragraph{PASS-008: Approved Input Promotion funktioniert}
17 model-promotierbare accepted Subjects wurden zu aktiven Approved Inputs
materialisiert; Open Questions wurden nicht erzwungen promoted.

\paragraph{PASS-009: Approved Engineering Information wird an Phase H übergeben}
17 Approved Subjects und 21 accepted semantic Relationships wurden gebunden.

\paragraph{PASS-010: Phase-H Readiness trennt Subject- und Relationship-Populationen}
Die Authority-Populationen werden getrennt geführt; sechs accepted Relationships
zu nicht model-promotierbaren Open Questions wurden als \texttt{Not projected}
beziehungsweise \texttt{intentionally\_not\_projected} behandelt.

\subsection{Finale Acceptance-Evidenz}

\paragraph{FINAL-G3: Project 000116}
Gate-3 Single-Source E2E: \texttt{PASS}. Reale externe SYSIDE-Validierung:
\texttt{PASS}. Human Publication Approval: \texttt{PASS}. Immutable Published
Output: \texttt{PASS}.

\paragraph{FINAL-MS: Project 308131}
Genuine-Multi-Source-Retest: \texttt{PASS}. BLK-002 final
\texttt{RESOLVED / LIVE RETESTED}.

\paragraph{FINAL-REG: finaler Integrationsstand}
Complete Repository Regression: \texttt{6335 passed, 12 skipped}. Diese Zahl
beschreibt die Breite der technischen Regressionsevidenz und ist kein
wissenschaftlicher Qualitäts- oder Korrektheitsscore.

\subsection{Verwendung des Registers im Haupttext}

Kapitel 7 verwendet nur thesis-relevante Beobachtungen aus diesem Register und
verweist bei Bedarf auf die zugehörigen Finding-IDs. Die Interpretation ihrer
Bedeutung erfolgt erst in Kapitel 8. UX- und Integrationsbeobachtungen werden nur
dort in den Haupttext übernommen, wo sie eine Forschungsfrage oder eine explizite
Systemgrenze unmittelbar stützen.
